% problem_id: S001-proposition-siblings-isomorphic
% source_id: S001 (Stacks Project)
% source_file: equiv.tex
% statement_locator: equiv.tex lines 2257--2267
% proof_locator: equiv.tex lines 2269--2457
% env: proposition
% status: text-extracted-verbatim (difficulty judgement pending, written by hand)
%
% The statement and proof below are the author's own text, extracted verbatim.
% Nothing is paraphrased, shortened, or supplemented.

\begin{proposition}
\label{proposition-siblings-isomorphic}
\begin{reference}
\cite[Proposition 2.16]{Orlov-K3}; the fact that we do not need
to assume vanishing of $\Ext^q(N, X)$ for $q > 0$ in the definition
of negative objects above is due to \cite{Canonaco-Stellari}.
\end{reference}
Let $F$ and $F'$ be siblings as in Definition \ref{definition-siblings}.
Assume that $F$ is fully faithful and that $\mathcal{A}$ has enough
negative objects (see above). Then $F$ and $F'$ are isomorphic functors.
\end{proposition}

\begin{proof}
By part (2) of Definition \ref{definition-siblings} the image of the functor
$F'$ is contained in the essential image of the functor $F$. Hence
the functor $H = F^{-1} \circ F'$ is a sibling of the identity functor.
This reduces us to the case described in the next paragraph.

\medskip\noindent
Let $\mathcal{D} = D^b(\mathcal{A})$. We have to show a sibling
$F : \mathcal{D} \to \mathcal{D}$ of the identity functor is
isomorphic to the identity functor. Given an object $X$ of $\mathcal{D}$
let us say $X$ has {\it width} $w = w(X)$ if $w \geq 0$ is minimal
such that there exists an integer $a \in \mathbf{Z}$ with $H^i(X) = 0$
for $i \not \in [a, a + w - 1]$. Since $F$ is a sibling of the identity
and since $F \circ [n] = [n] \circ F$ we are already given isomorphisms
$$
c_X : X \to F(X)
$$
for $w(X) \leq 1$ compatible with shifts. Moreover, if $X = A[-a]$ and
$X' = A'[-a]$ for some $A, A' \in \Ob(\mathcal{A})$ then for any morphism
$f : X \to X'$ the diagram
\begin{equation}
\label{equation-to-show}
\vcenter{
\xymatrix{
X \ar[d]_{c_X} \ar[r]_f &
X' \ar[d]^{c_{X'}} \\
F(X) \ar[r]^{F(f)} &
F(X')
}
}
\end{equation}
is commutative.

\medskip\noindent
Next, let us show that for any morphism $f : X \to X'$  with
$w(X), w(X') \leq 1$ the diagram (\ref{equation-to-show}) commutes.
If $X$ or $X'$ is zero, this is clear. If not then we can write
$X = A[-a]$ and $X' = A'[-a']$ for unique $A, A'$ in $\mathcal{A}$
and $a, a' \in \mathbf{Z}$. The case $a = a'$ was discussed above.
If $a' > a$, then $f = 0$ (Derived Categories, Lemma
\ref{derived-lemma-negative-exts}) and the result is clear.
If $a' < a$ then $f$ corresponds to an element
$\xi \in \Ext^q(A, A')$ with $q = a - a'$. Using Yoneda extensions, see
Derived Categories, Section \ref{derived-section-ext}, we can find
$A = A_0, A_1, \ldots, A_{q - 1}, A_q = A' \in \Ob(\mathcal{A})$ and
elements
$$
\xi_i \in \Ext^1(A_{i - 1}, A_i)
$$
such that $\xi$ is the composition $\xi_q \circ \ldots \circ \xi_1$.
In other words, setting $X_i = A_i[-a + i]$
we obtain morphisms
$$
X = X_0 \xrightarrow{f_1} X_1 \to \ldots \to X_{q - 1}
\xrightarrow{f_q} X_q = X'
$$
whose compostion is $f$. Since the commutativity of  (\ref{equation-to-show})
for $f_1, \ldots, f_q$ implies it for $f$, this reduces us to the case $q = 1$.
In this case after shifting we may assume we have a distinguished triangle
$$
A' \to E \to A \xrightarrow{f} A'[1]
$$
Observe that $E$ is an object of $\mathcal{A}$. Consider the following
diagram
$$
\xymatrix{
E \ar[d]_{c_E} \ar[r] &
A \ar[d]_{c_A} \ar[r]_f &
A'[1] \ar[d]^{c_{A'}[1]}
\ar@{..>}@<-1ex>[d]_\gamma \ar@{..>}[ld]^\epsilon \ar[r] &
E[1] \ar[d]^{c_E[1]} \\
F(E) \ar[r] &
F(A) \ar[r]^{F(f)} &
F(A')[1] \ar[r] &
F(E)[1]
}
$$
whose rows are distinguished triangles.
The square on the right commutes already but we don't yet know that
the middle square does. By the axioms of a triangulated category
we can find a morphism $\gamma$ which does make the diagram commute.
Then $\gamma - c_{A'}[1]$ composed with
$F(A')[1] \to F(E)[1]$ is zero hence we
can find $\epsilon : A'[1] \to F(A)$ such that
$\gamma - c_{A'}[1] = F(f) \circ \epsilon$. However, any arrow
$A'[1] \to F(A)$ is zero as it is a negative ext class
between objects of $\mathcal{A}$. Hence $\gamma = c_{A'}[1]$
and we conclude the middle square commutes too which is what we
wanted to show.

\medskip\noindent
To finish the proof we are going to argue by induction on $w$
that there exist isomorphisms $c_X : X \to F(X)$ for all
$X$ with $w(X) \leq w$ compatible with all morphisms between
such objects. The base case $w = 1$ was shown above. Assume
we know the result for some $w \geq 1$.

\medskip\noindent
Let $X$ be an object with $w(X) = w + 1$. Pick $a \in \mathbf{Z}$ with
$H^i(X) = 0$ for $i \not \in [a, a + w]$. Set $b = a + w$ so that
$H^b(X)$ is nonzero. Choose $N[-b] \to X$ as in Lemma \ref{lemma-good-map}.
Choose a distinguished diagram
$$
N[-b] \to X \to Y \to N[-b + 1]
$$
Computing the long exact cohomology sequence we find
$w(Y) \leq w$. Hence by induction we find the solid arrows
in the following diagram
$$
\xymatrix{
N[-b] \ar[r] \ar[d]_{c_N[-b]} &
X \ar[r] \ar@{..>}[d]_{c_{N[-b] \to X}} &
Y \ar[r] \ar[d]^{c_Y} &
N[-b + 1] \ar[d]^{c_N[-b + 1]} \\
F(N)[-b] \ar[r] &
F(X) \ar[r] &
F(Y) \ar[r] &
F(N)[-b + 1]
}
$$
We obtain the dotted arrow $c_{N[-b] \to X}$.
By Derived Categories, Lemma \ref{derived-lemma-uniqueness-third-arrow}
the dotted arrow is unique because $\Hom(X, F(N)[-b]) \cong \Hom(X, N[-b]) = 0$
by our choice of $N$. In fact, $c_{N[-b] \to X}$ is the unique dotted
arrow making the square with vertices $X, Y, F(X), F(Y)$ commute.

\medskip\noindent
Let $N'[-b] \to X$ be another map as in Lemma \ref{lemma-good-map}
and let us prove that $c_{N[-b] \to X} = c_{N'[-b] \to X}$.
Observe that the map $(N \oplus N')[-b] \to X$ also satisfies the
conditions of Lemma \ref{lemma-good-map}.
Thus we may assume $N'[-b] \to X$ factors
as $N'[-b] \to N[-b] \to X$ for some morphism $N' \to N$.
Choose distinguished triangles $N[-b] \to X \to Y \to N[-b + 1]$ and
$N'[-b] \to X \to Y' \to N'[-b + 1]$. By axiom TR3 we can find
a morphism $g : Y' \to Y$ which joint with $\text{id}_X$ and $N' \to N$
forms a morphism of triangles. Since we have
(\ref{equation-to-show}) for $g$ we conclude that
$$
(F(X) \to F(Y)) \circ c_{N'[-b] \to X} = (F(X) \to F(Y)) \circ c_{N[-b] \to X}
$$
The uniqueness of $c_{N[-b] \to X}$ pointed out in the construction
above now shows that $c_{N'[-b] \to X} = c_{N[-b] \to X}$.

\medskip\noindent
Thus we can now define for $X$ of width $w + 1$ the isomorphism
$c_X : X \to F(X)$ as the common value of the maps $c_{N[-b] \to X}$
where $N[-b] \to X$ is as in Lemma \ref{lemma-good-map}. To finish
the proof, we have to show that the diagrams (\ref{equation-to-show})
commute for all morphisms $f : X \to X'$ between objects with $w(X) \leq w + 1$
and $w(X') \leq w + 1$. Choose $a \leq b \leq a + w$ such that
$H^i(X) = 0$ for $i \not \in [a, b]$ and
$a' \leq b' \leq a' + w$ such that $H^i(X') = 0$ for
$i \not \in [a', b']$. We will use induction on
$(b' - a') + (b - a)$ to show the claim. (The base case
is when this number is zero which is OK because $w \geq 1$.)
We distinguish two cases.

\medskip\noindent
Case I: $b' < b$. In this case, by Lemma \ref{lemma-good-map-zero}
we may choose $N[-b] \to X$ as in Lemma \ref{lemma-good-map}
such that the composition $N[-b] \to X \to X'$ is zero.
Choose a distuiguished triangle $N[-b] \to X \to Y \to N[-b + 1]$. Since
$N[-b] \to X'$ is zero, we find that $f$ factors
as $X \to Y \to X'$. Since $H^i(Y)$ is nonzero only for $i \in [a, b - 1]$
we see by induction that (\ref{equation-to-show}) commutes for
$Y \to X'$. The diagram (\ref{equation-to-show}) commutes for
$X \to Y$ by construction if $w(X) = w + 1$ and by our first
induction hypothesis if $w(X) \leq w$.
Hence (\ref{equation-to-show}) commutes for $f$.

\medskip\noindent
Case II: $b' \geq b$. In this case we choose $N'[-b'] \to X'$
as in Lemma \ref{lemma-good-map}.
We may also assume that $\Hom(H^{b'}(X), N') = 0$ (this is
relevant only if $b' = b$), for example because we can
replace $N'$ by an object $N''$ which surjects onto $N' \oplus H^{b'}(X)$
and such that $\Hom(N' \oplus H^{b'}(X), N'') = 0$.
We choose a distinguished triangle
$N'[-b'] \to X' \to Y' \to N'[-b' + 1]$. Since
$\Hom(X, X') \to \Hom(X, Y')$ is injective by our choice of $N'$
(details omitted) the same is true for
$\Hom(X, F(X')) \to \Hom(X, F(Y'))$.
Hence it suffices in this case to check that
(\ref{equation-to-show}) commutes for the composition $X \to Y'$
of the morphisms $X \to X' \to Y'$.
Since $H^i(Y')$ is nonzero only for $i \in [a', b' - 1]$
we conclude by induction hypothesis.
\end{proof}
